\chapter{Обсуждение результатов}

\section{Сравнительный анализ решений и проверка подстановкой}

В результате работы программного комплекса получены численные решения системы уравнений варианта №~11 прямым методом Гаусса с выбором главного элемента и методом простой итерации:
\begin{equation}
    x_{\text{Gauss}} = \begin{pmatrix} 1{,}27343 \\ -2{,}33990 \\ -1{,}44510 \end{pmatrix},
    \qquad
    x_{\text{Iter}} = \begin{pmatrix} 1{,}27342 \\ -2{,}33989 \\ -1{,}44509 \end{pmatrix}.
\end{equation}
Максимальное расхождение между результатами методов Гаусса и простой итерации составляет:
\begin{multline}
    \|x_{\text{Gauss}} - x_{\text{Iter}}\|_\infty = \\
    = \max\big(|1{,}27343 - 1{,}27342|;\, |-2{,}33990 - (-2{,}33989)|; \\
    |-1{,}44510 - (-1{,}44509)|\big) = 10^{-5},
\end{multline}
что в точности соответствует заданной в условии точности $\varepsilon = 10^{-5}$.

\subsection{Проверка результатов подстановкой в исходную систему}

В соответствии с требованиями к отчету выполнена проверка полученных результатов путем непосредственной подстановки векторов $x_{\text{Gauss}}$ и $x_{\text{Iter}}$ в исходную систему уравнений~\eqref{eq:var11}.

Для решения методом Гаусса значения левых частей уравнений составляют:
\begin{equation}
    A x_{\text{Gauss}} = \begin{pmatrix}
        2{,}7 \cdot 1{,}27343 + 0{,}9 \cdot (-2{,}33990) - 1{,}5 \cdot (-1{,}44510) \\
        4{,}5 \cdot 1{,}27343 - 2{,}8 \cdot (-2{,}33990) + 6{,}7 \cdot (-1{,}44510) \\
        5{,}1 \cdot 1{,}27343 + 3{,}7 \cdot (-2{,}33990) - 1{,}4 \cdot (-1{,}44510)
    \end{pmatrix} = \begin{pmatrix}
        3{,}500001 \\
        2{,}599985 \\
        -0{,}139997
    \end{pmatrix}.
\end{equation}
Разности между левыми и правыми частями равны $|(Ax)_1 - b_1| = 1{,}0 \cdot 10^{-6}$, $|(Ax)_2 - b_2| = 1{,}5 \cdot 10^{-5}$, $|(Ax)_3 - b_3| = 3{,}0 \cdot 10^{-6}$.

Для решения методом простой итерации:
\begin{equation}
    A x_{\text{Iter}} = \begin{pmatrix}
        2{,}7 \cdot 1{,}27342 + 0{,}9 \cdot (-2{,}33989) - 1{,}5 \cdot (-1{,}44509) \\
        4{,}5 \cdot 1{,}27342 - 2{,}8 \cdot (-2{,}33989) + 6{,}7 \cdot (-1{,}44509) \\
        5{,}1 \cdot 1{,}27342 + 3{,}7 \cdot (-2{,}33989) - 1{,}4 \cdot (-1{,}44509)
    \end{pmatrix} = \begin{pmatrix}
        3{,}499968 \\
        2{,}599979 \\
        -0{,}140025
    \end{pmatrix}.
\end{equation}
Разности между левыми и правыми частями не превосходят $3{,}2 \cdot 10^{-5}$.

Полученные результаты практически точно совпадают с вектором свободных членов $b = (3{,}5;\, 2{,}6;\, -0{,}14)^T$, что полностью подтверждает правильность найденных решений. Погрешности для метода Гаусса вызваны округлением компонент вектора при выводе до пяти знаков, а для метода итераций согласуются с заданной точностью $\varepsilon = 10^{-5}$.

\section{Анализ сходимости итерационного процесса}

Скорость сходимости метода простой итерации определяется величиной нормы матрицы перехода $B$. В нашем случае кубическая норма матрицы составила:
\begin{equation}
    \|B\|_\infty = 0{,}892857 < 1.
\end{equation}

Поскольку норма $\|B\|_\infty$ меньше единицы, достаточное условие сходимости выполнено. При этом норма достаточно близка к единице ($0{,}893$), что обусловило плавную динамику уменьшения погрешности: для достижения требуемой точности $\varepsilon = 10^{-5}$ потребовалось 23 итерации. На последней итерации разность соседних приближений составила $\|x^{(23)} - x^{(22)}\|_\infty = 8{,}17 \cdot 10^{-6} < 10^{-5}$, что полностью удовлетворяет критерию окончания вычислений.

\section{Сравнение достоинств и недостатков методов}

Сравнительные характеристики прямого метода Гаусса с выбором главных элементов и метода простой итерации сведены в таблицу~\ref{tab:comparison}.

\begin{table}[H]
    \caption{Сравнительный анализ методов решения СЛАУ}
    \label{tab:comparison}
    \centering
    \small
    \begin{tabular}{|p{3.6cm}|p{5.4cm}|p{5.4cm}|}
        \hline
        \textbf{Характеристика} & \textbf{Метод Гаусса с выбором главных элементов} & \textbf{Метод простой итерации} \\
        \hline
        Класс метода & Прямой (точный) & Итерационный (приближенный) \\
        \hline
        Вычислительная сложность & $O(n^3)$ (для $n=3$: $\approx 18$ операций) & $O(k \cdot n^2)$ ($k=23$ итераций, $\approx 207$ операций) \\
        \hline
        Требования к матрице системы & $\det(A) \neq 0$ (невырожденность) & Наличие нормы $\|B\| < 1$ (диагональное преобладание) \\
        \hline
        Влияние ошибок машинного округления & Ошибки накапливаются в процессе прямого хода & Метод самоисправляемый; ошибки предыдущих шагов нивелируются \\
        \hline
        Необходимость начального приближения & Не требуется & Требуется вектор $x^{(0)}$ (обычно выбирается $c$) \\
        \hline
        Реакция на вырожденность / деление на 0 & Корректно определяет отсутствие решений или бесконечное число решений & При отсутствии диагонального преобладания процесс расходится \\
        \hline
        Область применения & Плотные матрицы умеренной размерности ($n \le 10^3$) & Разреженные матрицы сверхбольшой размерности ($n > 10^4$) \\
        \hline
    \end{tabular}
\end{table}

\section{Заключение}

В ходе выполнения лабораторной работы №~1 были успешно решены все поставленные задачи:
\begin{enumerate}
    \item Исходная система уравнений варианта №~11 исследована на совместность и единственность решения ($\det(A) = -66{,}321 \neq 0$, система совместна и определенна);
    \item Реализован прямой метод Гаусса с полным выбором главного элемента по матрице, исключающий деление на ноль и аварийное завершение программы;
    \item Реализован метод простой итерации с организацией хранения приближений в двух одномерных массивах;
    \item Исходная система успешно преобразована в эквивалентную систему со строгим диагональным преобладанием, обеспечив выполнение достаточного условия сходимости $\|B\|_\infty = 0{,}892857 < 1$;
    \item Проведена проверка найденных решений подстановкой в исходную систему уравнений (максимальные расхождения с вектором свободных членов составили $1{,}5 \cdot 10^{-5}$ и $3{,}2 \cdot 10^{-5}$ соответственно), что доказывает правильность и высокую точность вычислений.
\end{enumerate}
